\documentclass[10pt,a4paper]{article}
\usepackage[T1]{fontenc}
\usepackage{lmodern}
\usepackage{amsmath,amssymb,amsthm,mathtools}
\usepackage[margin=24mm]{geometry}
\usepackage{microtype}
\usepackage{booktabs,array,longtable}
\usepackage{xcolor}
\usepackage{hyperref}
\usepackage{fancyhdr}
\hypersetup{colorlinks=true,linkcolor=black,citecolor=black,urlcolor=blue,pdfauthor={Matthew J. Colbrook}}
\pagestyle{fancy}
\fancyhf{}
\fancyhead[L]{\small Partial result / MI-25}
\fancyhead[R]{\small 11 September 2026}
\fancyfoot[C]{\thepage}
\setlength{\headheight}{14pt}
\setlength{\parskip}{4pt}
\setlength{\parindent}{0pt}
\newtheorem{theorem}{Theorem}[section]
\newtheorem{proposition}[theorem]{Proposition}
\newtheorem{lemma}[theorem]{Lemma}
\newtheorem{corollary}[theorem]{Corollary}
\theoremstyle{remark}
\newtheorem{remark}[theorem]{Remark}
\newcommand{\M}{\mathbb M}
\newcommand{\C}{\mathbb C}
\newcommand{\R}{\mathbb R}
\newcommand{\tr}{\operatorname{tr}}
\newcommand{\diag}{\operatorname{diag}}
\newcommand{\spec}{\operatorname{spec}}
\newcommand{\rank}{\operatorname{rank}}
\newcommand{\abs}[1]{\lvert #1\rvert}
\newcommand{\norm}[1]{\lVert #1\rVert}
\newcommand{\opnorm}[1]{\lVert #1\rVert_{\infty}}
\newcommand{\schatten}[2]{\lVert #1\rVert_{#2}}
\newcommand{\symmod}{\mathcal S}
\newcommand{\maxmod}{\mathcal M}
\newcommand{\draftnotice}{\begin{center}\small\textit{Proposed mathematical resolution. Complete proof supplied; not submitted or peer reviewed.}\end{center}}

\usepackage{xurl}
\setlength{\emergencystretch}{3em}
\allowdisplaybreaks[1]

\title{MI-25: Partial Result}
\author{Matthew J. Colbrook\\
\small Department of Applied Mathematics and Theoretical Physics\\
\small University of Cambridge, Cambridge, United Kingdom\\
\small\href{mailto:m.colbrook@damtp.cam.ac.uk}{m.colbrook@damtp.cam.ac.uk}}
\date{11 September 2026}
\begin{document}
\maketitle
\textbf{Repository status: Partially resolved.}
The complete argument received an independent Codex-agent PASS review
for the stated partial result only.
The original draft was AI-assisted; the reviewed proof below is unchanged.
This is independent agent verification, not external human peer review or
formal proof certification. \href{https://github.com/ajt60gaibb/OpenProblemsInNLA/blob/main/references/colbrook-matrix-2026-09-11/verification/reviews/MI-25-review.md}{Detailed review and scope}.
\par\medskip
% BEGIN REVIEWED PROOF
\section{MI-25, endpoint only: the trace-norm Hlawka constant is infinite}
\label{sec:mi25}

For a norm $\|\cdot\|$, define the three-summand defect and the sum of pairwise defects by
\begin{align*}
 N(X_1,X_2,X_3)&=\sum_{j=1}^3\|X_j\|-\|X_1+X_2+X_3\|,\\
 D(X_1,X_2,X_3)&=\sum_{1\leq i<j\leq3}
 \bigl(\|X_i\|+\|X_j\|-\|X_i+X_j\|\bigr).
\end{align*}
The full repository problem asks for the optimal dimension-independent constants $C_p$ in $N\leq C_pD$ for Schatten norms. The result below resolves \emph{only the endpoint $p=1$}, which was left open in \cite[Problem 7]{AK2012}. It does not determine $C_p$ for $1<p<\infty$.

\begin{theorem}\label{thm:mi25}
No finite constant $C_1$ controls the three-summand trace-norm defect by the sum of pairwise trace-norm defects. This failure already occurs for real $2\times3$ rank-one matrices, or for real square $3\times3$ matrices after zero padding.
\end{theorem}

\begin{proof}
Let
\[
 u_1=\begin{pmatrix}1\\0\end{pmatrix},\qquad
 u_2=\begin{pmatrix}-1/2\\\sqrt3/2\end{pmatrix},\qquad
 u_3=\begin{pmatrix}-1/2\\-\sqrt3/2\end{pmatrix}.
\]
These unit vectors satisfy
\[
 \sum_ju_j=0,\qquad \sum_ju_ju_j^T=\frac32I_2,\qquad
 u_i^Tu_j=-\frac12\quad(i\ne j).
\]
For $0<t<1$, put
\[
 v_j=\begin{pmatrix}\sqrt{1-t}\,u_j\\\sqrt t\end{pmatrix}\in\R^3,
 \qquad X_j=u_jv_j^T\in\M_{2,3}(\R).
\]
Each $X_j$ has its only nonzero singular value equal to one. Furthermore,
\[
 \sum_jX_j=\begin{pmatrix}\frac32\sqrt{1-t}\,I_2&0\end{pmatrix},
 \qquad\left\|\sum_jX_j\right\|_1=3\sqrt{1-t}.
\]
For $i\ne j$, write $U=(u_i\ u_j)$ and $V=(v_i\ v_j)$. The two Gram matrices are
\[
 U^TU=\begin{pmatrix}1&-1/2\\-1/2&1\end{pmatrix},\qquad
 V^TV=\begin{pmatrix}1&(3t-1)/2\\(3t-1)/2&1\end{pmatrix}.
\]
They commute and are simultaneously diagonalized by $(1,1)^T$ and $(1,-1)^T$. The squared singular values of $UV^T=X_i+X_j$ are therefore
\[
 \frac{1+3t}{4},\qquad\frac{9(1-t)}4.
\]
It follows that every pair sum has trace norm
\[
 q(t)=\frac12\sqrt{1+3t}+\frac32\sqrt{1-t}.
\]
Consequently,
\[
 \frac{N(X_1,X_2,X_3)}{D(X_1,X_2,X_3)}
 =\frac{1-\sqrt{1-t}}{2-q(t)}.
\]
The denominator is positive for $t>0$: strict concavity of the square root gives
$\frac14\sqrt{1+3t}+\frac34\sqrt{1-t}<1$.
Expansion at zero gives
\[
 1-\sqrt{1-t}=\frac t2+O(t^2),\qquad
 2-q(t)=\frac34t^2+O(t^3).
\]
Thus the defect ratio is asymptotic to $2/(3t)$ and tends to $+\infty$.

Alternatively, the divergence can be checked without Taylor expansions. If
$g(t)=\sqrt{(1-t)(1+3t)}$, rationalization gives
\[
 2-q(t)=\frac{6t^2}{(1+t+g(t))(2+q(t))}
\]
and hence
\[
 \frac{1-\sqrt{1-t}}{2-q(t)}
 =\frac{(1+t+g(t))(2+q(t))}{6t(1+\sqrt{1-t})}
 \sim\frac{2}{3t}.
\]
Appending a zero row leaves every singular value and every defect unchanged, giving the square-matrix assertion.
\end{proof}

% END REVIEWED PROOF
\begin{thebibliography}{1}

\bibitem{AK2012}
Koenraad M.~R. Audenaert and Fuad Kittaneh.
\newblock Problems and conjectures in matrix and operator inequalities.
\newblock arXiv:1201.5232, 2012.
\newblock Version consulted: v3; Problem 5, printed page 3, and Problem 7,
  printed pages 12--13.
\newblock URL: \url{https://arxiv.org/abs/1201.5232}.

\end{thebibliography}

\end{document}
